% Solución explícita y clasificación del universo descendiente.

\section{Solution régulière explicite de la continuation cosmologique descendante}
\label{sec:ivv-descendant-explicit}

Considérons l'équation effective homogène d'Einstein--Cartan pour
\(n=3\),
\begin{equation}
 H^2=C\bigl(Aa^{-3}-Ba^{-6}\bigr),
 \qquad A,B,C>0.
 \label{eq:ivv-friedmann-descendant}
\end{equation}
Étant donnés \(t_b\in\mathbb R\), définissons
\begin{equation}
 a(t)^3=\frac BA+\frac{9AC}{4}(t-t_b)^2.
 \label{eq:ivv-scale-factor-descendant}
\end{equation}

\begin{theorem}[Existence non vide et complétude de la continuation descendante]
\label{thm:ivv-descendant-existence}
La fonction \eqref{eq:ivv-scale-factor-descendant} est une solution globale de
\eqref{eq:ivv-friedmann-descendant}. Elle possède un unique minimum
\(a_b=(B/A)^{1/3}>0\), satisfait
\[
 H(t_b)=0,
 \qquad
 \dot H(t_b)=\frac{3A^2C}{2B}>0,
\]
et définit sur \(\mathbb R\times\mathbb T^3\) une métrique FLRW lisse,
causalement orientable, avec des invariants de courbure bornés et des géodésiques
causales complètes. Les régions \(t<t_b\) et \(t>t_b\) réalisent la branche
parentale contractante et la continuation descendante expansive, unies par l'
hypersurface régulière \(t=t_b\).
\end{theorem}

\begin{proof}
En posant \(y=a^3\), on obtient
\[
 H=\frac{\dot y}{3y}=\frac{3AC(t-t_b)}{2y},
 \qquad
 C(Aa^{-3}-Ba^{-6})
 =C\frac{Ay-B}{y^2}
 =\frac{9A^2C^2(t-t_b)^2}{4y^2}=H^2.
\]
La positivité de \(B/A\) exclut \(a=0\). Aussi bien \(H\) que \(\dot H\)
sont bornés ; les invariants polynomiaux de courbure de la métrique FLRW
plate le sont également. Le temps cosmique parcourt \(\mathbb R\), et le paramètre
affine des géodésiques nulles contient \(\int a(t)\,dt\), qui diverge aux deux
extrémités parce que \(a(t)\asymp |t|^{2/3}\). La longueur propre temporelle diverge
également. Il n'existe donc pas de bord causal fini ni de singularité au
rebond.
\end{proof}

\begin{theorem}[Obstruction au cobordisme ramifié classique]
\label{thm:ivv-cobordism-obstruction}
Dans la classe des variétés lorentziennes lisses, non dégénérées, orientables dans
le temps et globalement hyperboliques, une évolution entre surfaces de
Cauchy compactes ne peut produire une ramification avec changement de topologie.
Tout feuilletage de Cauchy est difféomorphe à \(\mathbb R\times\Sigma\) ; par conséquent,
les feuillets initial et final sont difféomorphes. Le cobordisme ramifié fait l'objet d'une
obstruction au sein de cette classe.
\end{theorem}

Ainsi sont résolus les \(2\) objets géométriques : la continuation minimale
possède la solution complète précédente, et
la ramification classique avec changement topologique est exclue par le théorème d'
obstruction. Une géométrie dégénérée ou une théorie quantique de changement
topologique constitue un autre domaine et ne modifie pas ces résultats.
